% !TeX root = ../main.tex
\section{Methodology and Safety}\label{sec:method}
\rubric{Safety and zero-queue baselines}{30}{Safe setup, direct I/O and alignment, and an accurate QD=1
latency table (with \cref{sec:zero-queue}).}
\rubric{Reproducibility and data hygiene}{25}{Configs, raw data, scripts, repeated trials and
environment documentation.}

\subsection{Platform}
\newcommand{\acspinning}{a single I/O-issuing thread, not CPU-pinned (storage I/O is device-bound, not
CPU-bound, below the host-saturation queue depths of \cref{sec:queue-depth})}
\newcommand{\acsenvextra}{%
  SSD       & Micron 2210 \code{MTFDHBA512QFD}, \qty{512}{\giga\byte} QLC NVMe \cite{micron2210brief},
              firmware \code{P6MA001} (SMART log taken with root after the run, \cref{sec:anomalies}) \\
  Interface & PCIe \qty{8.0}{GT/s} $\times$4 (Gen3), current and maximum, from sysfs. Raw ceiling
              $8\,\text{GT/s}\times4\times\tfrac{128}{130}=\qty{3.94}{\giga\byte\per\second}$
              (\qty{3756}{\mebi\byte\per\second}; 128b/130b coding \cite{pcisig_pcie3}), roughly
              \qtyrange{3.2}{3.5}{\giga\byte\per\second} after TLP/DLLP overhead \\
  Block layer & scheduler \code{none}, \code{nr_requests} 1023, 16 I/O queues, volatile write cache
              \enquote{write back} \\
  Target    & disposable files on the SSD's ext4 \code{/} partition (\code{nvme0n1p3}, \code{relatime},
              no online discard) via \code{ACS_SCRATCH}: a \qty{2}{\gibi\byte} test file for the sweeps
              and a \qty{60}{\gibi\byte} file for the device-state runs \\
  Temperature & NVMe composite sensor (\code{hwmon*/temp1_input}), logged per repetition and every
              second in the sustained runs \\
}
%% Measurement environment table (common requirement: measurement discipline).
%% Values come from common/system_info.tex. Before \input-ing this file a project may define
%%   \acspinning   how runs were pinned/isolated
%%   \acsenvextra  extra rows, e.g.  \newcommand{\acsenvextra}{SSD & model, firmware, PCIe link \\}
\providecommand{\acspinning}{\todo{e.g.\ \code{taskset -c 2} on a P-core; isolcpus if used}}
\providecommand{\acsenvextra}{}
\begin{table}[tbp]
  \centering\small
  \caption{Measurement environment.}\label{tab:environment}
  \begin{tabularx}{\linewidth}{@{}>{\raggedright\arraybackslash}p{0.2\linewidth}X@{}}
    \toprule
    CPU               & \sysCPU{} (\sysArch) \\
    SIMD ISA          & \sysISA \\
    Topology          & \sysCores; NUMA: \sysNUMA \\
    SMT               & \sysSMT \\
    Caches            & \sysCaches \\
    Memory            & \sysMemory; DIMMs: \sysDIMM \\
    Pages             & base \sysPageSize; THP: \sysHugePages \\
    OS / kernel       & \sysOS; \sysKernel \\
    Frequency policy  & governor: \sysGovernor; driver: \sysDriver; turbo: \sysTurbo; EPP: \sysEPP \\
    Pinning / isolation & \acspinning; \sysIsolation \\
    Compilers         & \sysGCC; \sysClang \\
    Tools             & perf: \sysPerf{} (perf\_event\_paranoid \sysPerfParanoid); FIO: \sysFio; MLC: \sysMLC \\
    GPU (optional)    & \sysGPU; CUDA \sysCUDA \\
    \acsenvextra
    Env.\ snapshot     & \sysCollected{} (each raw run directory also stores its own \code{env.txt}) \\
    \bottomrule
  \end{tabularx}
\end{table}


The machine also has a SATA HDD (\code{/dev/sda}); each run's \code{env.txt} resolves the scratch
partition to its parent disk and records its \code{rotational} flag, model, firmware, link and
temperature, so a run cannot be silently mislabelled (\cref{sec:ai-use} describes how this check once
failed).

\subsection{Data-loss safeguards}
\begin{itemize}
  \item \textbf{Targets.} Only disposable files in a dedicated scratch directory are written, and both are
        deleted when \code{run.sh} exits, including on interrupt.
  \item \textbf{Gate.} \code{scripts/safety.sh} runs before every experiment and refuses block or character
        devices (also via symlinks), mountpoints, directories, paths outside the scratch directory and
        targets without enough free space; the scratch directory itself may not be \code{/},
        \code{\$HOME}, a filesystem root or foreign-owned. It only \code{stat}s paths
        (\cref{sec:validation}).
  \item \textbf{Benchmark.} \code{blkbench} \code{lstat}s its target, refuses anything but a regular file
        and opens with \code{O_NOFOLLOW}; no raw device is ever opened.
  \item \textbf{Space.} The large device-state file is refused unless \qty{10}{\gibi\byte} stay free.
\end{itemize}

\subsection{I/O engines and configuration}
The primary engine, \code{build/blkbench}, is a small O\_DIRECT \code{libaio} benchmark (the engine FIO
also uses \cite{axboe_fio}). It keeps exactly QD requests in flight in a closed loop and records every
I/O's submit-to-completion latency. FIO 3.42 runs the same points through \code{fio/job.fio}
(\cref{sec:validation}), and \code{ENGINE=fio} routes every sweep through FIO.
\begin{knobs}
  Engine          & \code{libaio}; \code{O_DIRECT}; 4\,KiB-aligned buffers and offsets; 64-bit uniform
                    random offsets (splitmix64) \\
  Timing          & \qty{4}{\second} per repetition (\qty{6}{\second} for tail); write points also stop after
                    \qty{256}{\mebi\byte}; preconditioning absorbs warm-up \\
  Repetitions     & 5 per point (3 for block-size writes; 3 per pass, 6 in total, for the working set);
                    median with inter-quartile-range (IQR) bars \\
  Data pattern    & incompressible xorshift64* buffers, fixed per queue slot (we did not test whether the
                    2210 deduplicates) \\
  Latency metric  & per-I/O submit$\to$completion; p50/p95/p99/p99.9 from the full sorted sample \\
  Preconditioning & sequential O\_DIRECT write over the test file at the start of every experiment and
                    before every sequential-read group (\cref{sec:layout}) \\
  Driver          & \code{SSD_Performance_Profiling/scripts/run.sh}; \code{src/blkbench.c} \\
\end{knobs}

\subsection{Execution, run order and idle periods}
All results except \code{pitfall} and \code{layout} come from one scripted run,
\code{scripts/rerun_clean.sh} (run \code{20261022-021116}, 2026-10-22, 02:11--04:08 EDT;
\cref{tab:budget}), started by the author in a terminal. Beforehand we set the governor to
\code{performance}, disabled turbo, SMT and ACPI GPE~0x42, ran \code{sudo fstrim -v /}
(\qty{42.8}{\gibi\byte} trimmed, about 01:41) and left the machine idle for \qty{30}{\minute}. The script
then checked free space (\qty{77}{\gibi\byte}), the TRIM in the system journal and the drive itself:
\qty{0}{\mebi\byte} read or written and \qty{0}{\percent} busy over \qty{10}{\second}.

The run puts the large sequential fill first, on the clean drive, and gives every later experiment an
idle period: \qty{20}{\minute} before \code{working_set}, \qty{10}{\minute} before \code{tail} and
\qty{5}{\minute} before each other step. Write points also wait before they start: \qty{120}{\second} in
\code{zero_queue} and \code{rw_mix}, \qty{15}{\second} in \code{block_size} and \code{queue_depth}.
\code{pitfall} and \code{layout} were not repeated; they come from our first run
(\code{20261021-103201}).

\textbf{Idle period.} An \emph{idle period} means that the drive received no I/O from our scripts for a
stated time (\qtyrange{5}{30}{\minute}), so that the controller could finish background SLC-to-QLC folding
and garbage collection; an \emph{idle drive} is one measured after such a period.

\paragraph{Superseded runs.} Two earlier runs are not used for the results. The first
(\code{run.sh all}, then \code{SUSTAIN_GIB=60 run.sh device_state}; runs \code{20261021-103201} and
\code{20261021-105439}, 2026-10-21) ran all experiments back to back without TRIM or idle periods, so each
inherited the SLC-cache fill and background work of the one before. A partial \enquote{lean} re-run that
evening (\code{20261021-220614}) re-measured only the most affected points, with a short fill, and could
only be merged with the first. Because of these run-order confounds both were replaced by the run above;
their raw data stay in the repository, and we cite them only where they show how drive state changes a
result.

\begin{table}[!htbp]
  \centering\small
  \caption{The reported run (\code{rerun_clean.sh}, 2026-10-22) in execution order: wall-clock time of
  each step from the run log, including its idle period, and the most it could write (byte caps plus
  preconditioning). \code{pitfall} and \code{layout} are from the first run.}\label{tab:budget}
  \begin{tabularx}{\linewidth}{@{}lrrX@{}}
    \toprule
    Step & Time & Max written & Notes \\
    \midrule
    \code{sustained}   & \qty{16.2}{\minute} & \qty{72}{\gibi\byte} & \qty{60}{\gibi\byte} fill $+$ $\le$\qty{8}{\gibi\byte} random $+$ \qty{60}{\second} idle $+$ $\le$\qty{4}{\gibi\byte} recovery \\
    \code{working_set} & \qty{24.1}{\minute} & \qty{6}{\gibi\byte}  & \qty{20}{\minute} idle; on the \qty{60}{\gibi\byte} file, which it then deletes \\
    \code{tail}        & \qty{12.6}{\minute} & \qty{2}{\gibi\byte}  & \qty{10}{\minute} idle; fresh \qty{2}{\gibi\byte} precondition; reads only \\
    \code{zero_queue}  & \qty{9.9}{\minute}  & \qty{7}{\gibi\byte}  & \qty{5}{\minute} idle; \qty{120}{\second} before each write point \\
    \code{rw_mix}      & \qty{11.6}{\minute} & \qty{4}{\gibi\byte}  & \qty{5}{\minute} idle; \qty{120}{\second} before each mix with writes \\
    \code{block_size}  & \qty{22.6}{\minute} & \qty{24}{\gibi\byte} & \qty{5}{\minute} idle; 56 points; \qty{15}{\second} before each write point \\
    \code{queue_depth} & \qty{11.0}{\minute} & \qty{14}{\gibi\byte} & \qty{5}{\minute} idle; reads and writes, QD 1--256 \\
    \code{fio_check}   & \qty{9.0}{\minute}  & \qty{10}{\gibi\byte} & \qty{5}{\minute} idle; 7 points, both engines, order alternating \\
    \midrule
    \code{pitfall}     & \qty{0.8}{\minute}  & \qty{2}{\gibi\byte}  & first run; reads only (plus re-fill) \\
    \code{layout}      & \qty{2.2}{\minute}  & \qty{6}{\gibi\byte}  & first run; includes a \qty{4}{\gibi\byte} random overwrite \\
    \bottomrule
  \end{tabularx}
\end{table}
